%%%%%%%%%%%%%%%%%%%%%%%%%%%%%%%%%%%%%%%%%%%%%%%%%%%%%%%%%%%%
% 4.6 KNOT THEORY
% The Composition of Advanced Mathematics
%%%%%%%%%%%%%%%%%%%%%%%%%%%%%%%%%%%%%%%%%%%%%%%%%%%%%%%%%%%%

\section{Knot Theory}
\label{sec:knot-theory}

\prerequisite{
Point-Set Topology, Algebraic Topology, Geometric Topology,
basic Group Theory, and introductory Manifold Theory.
}

\learningobjectives{
By the end of this section, the reader should be able to:
\begin{itemize}
    \item define knots and links mathematically;
    \item distinguish a knot from a knot diagram;
    \item explain ambient isotopy and knot equivalence;
    \item identify the unknot, trefoil, figure-eight knot, and torus knots;
    \item describe oriented knots and mirror images;
    \item explain regular projections and crossings;
    \item use Reidemeister moves to describe diagram equivalence;
    \item define crossing number;
    \item define connected sums of knots;
    \item distinguish prime and composite knots;
    \item define linking number;
    \item explain knot complements and knot groups;
    \item compute elementary knot-group presentations from diagrams;
    \item state the Wirtinger presentation;
    \item define knot colorings and explain their use as invariants;
    \item define the Alexander polynomial conceptually;
    \item define the Jones polynomial conceptually;
    \item explain the role of the HOMFLY-PT polynomial;
    \item define Seifert surfaces and knot genus;
    \item describe the Seifert algorithm;
    \item explain braids and braid groups;
    \item state Alexander's Theorem for closed braids;
    \item describe Markov moves conceptually;
    \item explain torus knots;
    \item distinguish knots, links, and spatial graphs;
    \item explain hyperbolic knot complements;
    \item describe Dehn surgery on knots;
    \item introduce knot concordance and slice knots;
    \item explain applications of knot theory.
\end{itemize}
}

%%%%%%%%%%%%%%%%%%%%%%%%%%%%%%%%%%%%%%%%%%%%%%%%%%%%%%%%%%%%
\subsection{What Is Knot Theory?}
\label{subsec:what-is-knot-theory}
%%%%%%%%%%%%%%%%%%%%%%%%%%%%%%%%%%%%%%%%%%%%%%%%%%%%%%%%%%%%

Knot theory studies embeddings of circles in three-dimensional space.

Informally, imagine tying a knot in a piece of string and then joining the
two ends together.

Once the ends are joined, the knot cannot be untied by pulling an end
through the knot.

Mathematically, the string becomes a copy of the circle

\[
S^1,
\]

and the surrounding space is usually taken to be

\[
S^3
\]

or equivalently \(\R^3\) together with a point at infinity.

Thus the fundamental object is

\[
\boxed{
S^1\hookrightarrow S^3.
}
\]

\begin{conceptbox}
Knot theory studies how circles can be embedded in three-dimensional space
and determines when two such embeddings can be continuously deformed into
one another without cutting or passing through themselves.
\end{conceptbox}

%%%%%%%%%%%%%%%%%%%%%%%%%%%%%%%%%%%%%%%%%%%%%%%%%%%%%%%%%%%%
\subsection{Knots}
\label{subsec:mathematical-knots}
%%%%%%%%%%%%%%%%%%%%%%%%%%%%%%%%%%%%%%%%%%%%%%%%%%%%%%%%%%%%

\begin{definitionbox}{Knot}
A \emph{knot} is an embedding

\[
\boxed{
K:S^1\hookrightarrow S^3.
}
\]

The image

\[
K(S^1)
\]

is also commonly called the knot.
\end{definitionbox}

Depending on context, knots may be considered in the smooth, piecewise-linear,
or tame topological category.

Classical knot theory generally excludes wild embeddings unless explicitly
stated otherwise.

%%%%%%%%%%%%%%%%%%%%%%%%%%%%%%%%%%%%%%%%%%%%%%%%%%%%%%%%%%%%
\subsection{Why Use the Three-Sphere?}
\label{subsec:why-three-sphere-knots}
%%%%%%%%%%%%%%%%%%%%%%%%%%%%%%%%%%%%%%%%%%%%%%%%%%%%%%%%%%%%

Classical knots are often drawn in

\[
\R^3.
\]

However,

\[
\boxed{
S^3\cong\R^3\cup\{\infty\}.
}
\]

Therefore a knot in \(\R^3\) can naturally be viewed as a knot in \(S^3\).

The compact space \(S^3\) is frequently more convenient for topology.

%%%%%%%%%%%%%%%%%%%%%%%%%%%%%%%%%%%%%%%%%%%%%%%%%%%%%%%%%%%%
\subsection{Knot Equivalence}
\label{subsec:knot-equivalence}
%%%%%%%%%%%%%%%%%%%%%%%%%%%%%%%%%%%%%%%%%%%%%%%%%%%%%%%%%%%%

Two knots are considered equivalent when one can be continuously moved into
the other through the surrounding space without cutting the knot or allowing
it to pass through itself.

The appropriate equivalence relation is ambient isotopy.

\begin{definitionbox}{Ambient Isotopy of Knots}
Two knots \(K_0\) and \(K_1\) are ambient isotopic if there exists a
continuous family of homeomorphisms

\[
\boxed{
F_t:S^3\to S^3,
\qquad
0\leq t\leq1,
}
\]

such that

\[
F_0=\operatorname{id}_{S^3}
\]

and

\[
\boxed{
F_1(K_0)=K_1.
}
\]
\end{definitionbox}

%%%%%%%%%%%%%%%%%%%%%%%%%%%%%%%%%%%%%%%%%%%%%%%%%%%%%%%%%%%%
\subsection{The Unknot}
\label{subsec:unknot}
%%%%%%%%%%%%%%%%%%%%%%%%%%%%%%%%%%%%%%%%%%%%%%%%%%%%%%%%%%%%

The simplest knot is the \emph{unknot}.

It is represented by an ordinary round circle:

\[
\boxed{
K_0\cong S^1.
}
\]

Every knot diagram with no crossings represents the unknot.

A major recurring problem is deciding whether a complicated diagram also
represents the unknot.

%%%%%%%%%%%%%%%%%%%%%%%%%%%%%%%%%%%%%%%%%%%%%%%%%%%%%%%%%%%%
\subsection{Nontrivial Knots}
\label{subsec:nontrivial-knots}
%%%%%%%%%%%%%%%%%%%%%%%%%%%%%%%%%%%%%%%%%%%%%%%%%%%%%%%%%%%%

A knot that is not equivalent to the unknot is called nontrivial.

Two famous examples are

\[
\boxed{
\text{trefoil knot}
}
\]

and

\[
\boxed{
\text{figure-eight knot}.
}
\]

The trefoil has minimum crossing number \(3\), while the figure-eight knot
has minimum crossing number \(4\).

%%%%%%%%%%%%%%%%%%%%%%%%%%%%%%%%%%%%%%%%%%%%%%%%%%%%%%%%%%%%
\subsection{Knot Diagrams}
\label{subsec:knot-diagrams}
%%%%%%%%%%%%%%%%%%%%%%%%%%%%%%%%%%%%%%%%%%%%%%%%%%%%%%%%%%%%

A knot is three-dimensional, but it is usually represented by a
two-dimensional projection.

At each crossing, one records which strand passes over the other.

A knot diagram therefore contains two kinds of information:

\[
\boxed{
\text{planar projection}
+
\text{over/under crossing data}.
}
\]

\begin{warningbox}
A knot diagram is not the knot itself.

Different diagrams may represent the same knot.
\end{warningbox}

%%%%%%%%%%%%%%%%%%%%%%%%%%%%%%%%%%%%%%%%%%%%%%%%%%%%%%%%%%%%
\subsection{Regular Projections}
\label{subsec:regular-knot-projections}
%%%%%%%%%%%%%%%%%%%%%%%%%%%%%%%%%%%%%%%%%%%%%%%%%%%%%%%%%%%%

A regular projection is chosen so that the diagram has only ordinary double
crossings.

One avoids degeneracies such as

\[
\boxed{
\text{triple crossings},
\quad
\text{tangencies},
\quad
\text{vertices at crossings}.
}
\]

After a small perturbation, tame knots admit regular projections.

%%%%%%%%%%%%%%%%%%%%%%%%%%%%%%%%%%%%%%%%%%%%%%%%%%%%%%%%%%%%
\subsection{Crossings}
\label{subsec:knot-crossings}
%%%%%%%%%%%%%%%%%%%%%%%%%%%%%%%%%%%%%%%%%%%%%%%%%%%%%%%%%%%%

At a crossing, one arc passes over another.

Diagrammatically, the under-strand is interrupted while the over-strand is
drawn continuously.

The over/under information is essential.

If it is removed, one generally cannot reconstruct the knot.

%%%%%%%%%%%%%%%%%%%%%%%%%%%%%%%%%%%%%%%%%%%%%%%%%%%%%%%%%%%%
\subsection{Reidemeister Moves}
\label{subsec:reidemeister-moves}
%%%%%%%%%%%%%%%%%%%%%%%%%%%%%%%%%%%%%%%%%%%%%%%%%%%%%%%%%%%%

There are three fundamental local moves on knot diagrams.

They are called the Reidemeister moves.

\[
\boxed{
R_1,\qquad R_2,\qquad R_3.
}
\]

They correspond to local changes that do not alter the knot type.

%%%%%%%%%%%%%%%%%%%%%%%%%%%%%%%%%%%%%%%%%%%%%%%%%%%%%%%%%%%%
\subsection{Reidemeister Move I}
\label{subsec:reidemeister-one}
%%%%%%%%%%%%%%%%%%%%%%%%%%%%%%%%%%%%%%%%%%%%%%%%%%%%%%%%%%%%

The first Reidemeister move adds or removes a twist.

Schematically,

\[
\boxed{
\text{straight arc}
\longleftrightarrow
\text{arc with one loop}.
}
\]

This changes the number of crossings by one.

%%%%%%%%%%%%%%%%%%%%%%%%%%%%%%%%%%%%%%%%%%%%%%%%%%%%%%%%%%%%
\subsection{Reidemeister Move II}
\label{subsec:reidemeister-two}
%%%%%%%%%%%%%%%%%%%%%%%%%%%%%%%%%%%%%%%%%%%%%%%%%%%%%%%%%%%%

The second Reidemeister move adds or removes two crossings.

It corresponds to moving one strand across another without changing the
underlying knot type.

Thus

\[
\boxed{
\Delta c=\pm2,
}
\]

where \(c\) denotes the number of crossings in the diagram.

%%%%%%%%%%%%%%%%%%%%%%%%%%%%%%%%%%%%%%%%%%%%%%%%%%%%%%%%%%%%
\subsection{Reidemeister Move III}
\label{subsec:reidemeister-three}
%%%%%%%%%%%%%%%%%%%%%%%%%%%%%%%%%%%%%%%%%%%%%%%%%%%%%%%%%%%%

The third Reidemeister move slides one strand across a crossing of two other
strands.

It preserves the number of crossings.

This move is especially important when proving that algebraic expressions
constructed from diagrams are knot invariants.

%%%%%%%%%%%%%%%%%%%%%%%%%%%%%%%%%%%%%%%%%%%%%%%%%%%%%%%%%%%%
\subsection{Reidemeister Theorem}
\label{subsec:reidemeister-theorem}
%%%%%%%%%%%%%%%%%%%%%%%%%%%%%%%%%%%%%%%%%%%%%%%%%%%%%%%%%%%%

\begin{theorem}[Reidemeister Theorem]
Two tame knot diagrams represent equivalent knots if and only if one diagram
can be transformed into the other by a finite sequence of Reidemeister moves
and planar isotopies.
\end{theorem}

Thus a three-dimensional topological equivalence problem becomes a
two-dimensional diagrammatic problem.

%%%%%%%%%%%%%%%%%%%%%%%%%%%%%%%%%%%%%%%%%%%%%%%%%%%%%%%%%%%%
\subsection{Knot Invariants}
\label{subsec:knot-invariants}
%%%%%%%%%%%%%%%%%%%%%%%%%%%%%%%%%%%%%%%%%%%%%%%%%%%%%%%%%%%%

A knot invariant assigns the same value to equivalent knots.

If

\[
K_1\cong K_2,
\]

then an invariant \(I\) must satisfy

\[
\boxed{
I(K_1)=I(K_2).
}
\]

Therefore,

\[
\boxed{
I(K_1)\neq I(K_2)
\Longrightarrow
K_1\not\cong K_2.
}
\]

The converse need not hold.

Two different knots may have the same value of a particular invariant.

%%%%%%%%%%%%%%%%%%%%%%%%%%%%%%%%%%%%%%%%%%%%%%%%%%%%%%%%%%%%
\subsection{Crossing Number}
\label{subsec:crossing-number}
%%%%%%%%%%%%%%%%%%%%%%%%%%%%%%%%%%%%%%%%%%%%%%%%%%%%%%%%%%%%

\begin{definitionbox}{Crossing Number}
The \emph{crossing number} of a knot \(K\), denoted

\[
\boxed{
c(K),
}
\]

is the smallest number of crossings among all diagrams representing \(K\).
\end{definitionbox}

For the unknot,

\[
\boxed{
c(K_0)=0.
}
\]

For the trefoil,

\[
\boxed{
c(3_1)=3.
}
\]

For the figure-eight knot,

\[
\boxed{
c(4_1)=4.
}
\]

%%%%%%%%%%%%%%%%%%%%%%%%%%%%%%%%%%%%%%%%%%%%%%%%%%%%%%%%%%%%
\subsection{Knot Tables}
\label{subsec:knot-tables}
%%%%%%%%%%%%%%%%%%%%%%%%%%%%%%%%%%%%%%%%%%%%%%%%%%%%%%%%%%%%

Knots are commonly indexed by crossing number.

For example,

\[
3_1
\]

denotes the first prime knot with crossing number \(3\).

Likewise,

\[
4_1
\]

denotes the first prime knot with crossing number \(4\).

The subscript distinguishes knots having the same crossing number.

%%%%%%%%%%%%%%%%%%%%%%%%%%%%%%%%%%%%%%%%%%%%%%%%%%%%%%%%%%%%
\subsection{Oriented Knots}
\label{subsec:oriented-knots}
%%%%%%%%%%%%%%%%%%%%%%%%%%%%%%%%%%%%%%%%%%%%%%%%%%%%%%%%%%%%

A knot can be assigned a direction around its circle.

Such a knot is called an \emph{oriented knot}.

An orientation is represented by arrows on the knot diagram.

Orientation becomes important for invariants such as

\[
\boxed{
\text{linking number}
}
\]

and for algebraic constructions involving crossings.

%%%%%%%%%%%%%%%%%%%%%%%%%%%%%%%%%%%%%%%%%%%%%%%%%%%%%%%%%%%%
\subsection{Mirror Images}
\label{subsec:knot-mirror-images}
%%%%%%%%%%%%%%%%%%%%%%%%%%%%%%%%%%%%%%%%%%%%%%%%%%%%%%%%%%%%

The mirror image of a knot \(K\) is denoted by notation such as

\[
\boxed{
\overline{K}.
}
\]

It is obtained by reflecting the knot in a mirror.

In a diagram, this can be represented by switching every crossing.

Some knots are equivalent to their mirror images.

Others are not.

%%%%%%%%%%%%%%%%%%%%%%%%%%%%%%%%%%%%%%%%%%%%%%%%%%%%%%%%%%%%
\subsection{Chirality}
\label{subsec:knot-chirality}
%%%%%%%%%%%%%%%%%%%%%%%%%%%%%%%%%%%%%%%%%%%%%%%%%%%%%%%%%%%%

A knot equivalent to its mirror image is called \emph{amphichiral}.

A knot not equivalent to its mirror image is called \emph{chiral}.

The trefoil is chiral.

The figure-eight knot is amphichiral.

Thus knot theory can distinguish a form of mathematical handedness.

%%%%%%%%%%%%%%%%%%%%%%%%%%%%%%%%%%%%%%%%%%%%%%%%%%%%%%%%%%%%
\subsection{Connected Sum of Knots}
\label{subsec:connected-sum-knots}
%%%%%%%%%%%%%%%%%%%%%%%%%%%%%%%%%%%%%%%%%%%%%%%%%%%%%%%%%%%%

Two knots can be combined to form a new knot.

Remove a small arc from each knot and join the resulting endpoints in the
standard way.

The resulting knot is the connected sum

\[
\boxed{
K_1\#K_2.
}
\]

For oriented knots, connected sum is well defined up to knot equivalence.

%%%%%%%%%%%%%%%%%%%%%%%%%%%%%%%%%%%%%%%%%%%%%%%%%%%%%%%%%%%%
\subsection{Prime Knots}
\label{subsec:prime-knots}
%%%%%%%%%%%%%%%%%%%%%%%%%%%%%%%%%%%%%%%%%%%%%%%%%%%%%%%%%%%%

\begin{definitionbox}{Prime Knot}
A nontrivial knot \(K\) is \emph{prime} if

\[
K=K_1\#K_2
\]

implies that at least one of \(K_1\) or \(K_2\) is the unknot.
\end{definitionbox}

A knot that is a connected sum of two nontrivial knots is called composite.

%%%%%%%%%%%%%%%%%%%%%%%%%%%%%%%%%%%%%%%%%%%%%%%%%%%%%%%%%%%%
\subsection{Prime Decomposition of Knots}
\label{subsec:prime-decomposition-knots}
%%%%%%%%%%%%%%%%%%%%%%%%%%%%%%%%%%%%%%%%%%%%%%%%%%%%%%%%%%%%

Every nontrivial tame knot can be expressed as a connected sum of prime
knots.

Moreover, this decomposition is unique up to order and knot equivalence.

Thus knot theory has a factorization theorem analogous to

\[
\boxed{
\text{integer prime factorization}.
}
\]

%%%%%%%%%%%%%%%%%%%%%%%%%%%%%%%%%%%%%%%%%%%%%%%%%%%%%%%%%%%%
\subsection{Links}
\label{subsec:links}
%%%%%%%%%%%%%%%%%%%%%%%%%%%%%%%%%%%%%%%%%%%%%%%%%%%%%%%%%%%%

A link generalizes a knot by allowing several circles.

\begin{definitionbox}{Link}
An \(m\)-component link is an embedding

\[
\boxed{
L:
\bigsqcup_{i=1}^{m}S^1
\hookrightarrow
S^3.
}
\]
\end{definitionbox}

A knot is therefore a one-component link.

%%%%%%%%%%%%%%%%%%%%%%%%%%%%%%%%%%%%%%%%%%%%%%%%%%%%%%%%%%%%
\subsection{The Unlink}
\label{subsec:unlink}
%%%%%%%%%%%%%%%%%%%%%%%%%%%%%%%%%%%%%%%%%%%%%%%%%%%%%%%%%%%%

The simplest \(m\)-component link is the unlink.

Its components can be separated into disjoint balls so that each component
is an unknot.

For two components, this looks like two unlinked circles.

%%%%%%%%%%%%%%%%%%%%%%%%%%%%%%%%%%%%%%%%%%%%%%%%%%%%%%%%%%%%
\subsection{The Hopf Link}
\label{subsec:hopf-link}
%%%%%%%%%%%%%%%%%%%%%%%%%%%%%%%%%%%%%%%%%%%%%%%%%%%%%%%%%%%%

The Hopf link is the simplest nontrivial two-component link.

Its two components are linked once.

For a suitable orientation,

\[
\boxed{
|\operatorname{lk}(L_1,L_2)|=1.
}
\]

It provides the basic example for the linking number.

%%%%%%%%%%%%%%%%%%%%%%%%%%%%%%%%%%%%%%%%%%%%%%%%%%%%%%%%%%%%
\subsection{Crossing Signs}
\label{subsec:crossing-signs}
%%%%%%%%%%%%%%%%%%%%%%%%%%%%%%%%%%%%%%%%%%%%%%%%%%%%%%%%%%%%

For an oriented link diagram, each crossing between two oriented strands can
be assigned a sign

\[
\boxed{
+1
\qquad\text{or}\qquad
-1.
}
\]

The sign depends on the relative orientations of the over-strand,
under-strand, and the orientation of the plane.

Crossing signs are used in linking number and writhe.

%%%%%%%%%%%%%%%%%%%%%%%%%%%%%%%%%%%%%%%%%%%%%%%%%%%%%%%%%%%%
\subsection{Linking Number}
\label{subsec:linking-number}
%%%%%%%%%%%%%%%%%%%%%%%%%%%%%%%%%%%%%%%%%%%%%%%%%%%%%%%%%%%%

Let \(L_1\) and \(L_2\) be two distinct oriented components of a link.

The linking number is

\[
\boxed{
\operatorname{lk}(L_1,L_2)
=
\frac12
\sum_{c}
\operatorname{sgn}(c),
}
\]

where the sum is taken only over crossings involving one strand from each
component.

The linking number is invariant under ambient isotopy of oriented links.

%%%%%%%%%%%%%%%%%%%%%%%%%%%%%%%%%%%%%%%%%%%%%%%%%%%%%%%%%%%%
\subsection{Worked Example: Linking Number}
\label{subsec:worked-linking-number}
%%%%%%%%%%%%%%%%%%%%%%%%%%%%%%%%%%%%%%%%%%%%%%%%%%%%%%%%%%%%

\begin{workedexamplebox}{Computing a Linking Number}

Suppose two oriented link components have four mutual crossings with signs

\[
+1,\quad+1,\quad+1,\quad-1.
\]

Then

\[
\sum_c\operatorname{sgn}(c)
=
1+1+1-1
=
2.
\]

Therefore

\[
\operatorname{lk}(L_1,L_2)
=
\frac12(2).
\]

\begin{answerbox}
\[
\boxed{
\operatorname{lk}(L_1,L_2)=1.
}
\]
\end{answerbox}
\end{workedexamplebox}

%%%%%%%%%%%%%%%%%%%%%%%%%%%%%%%%%%%%%%%%%%%%%%%%%%%%%%%%%%%%
\subsection{Limitations of Linking Number}
\label{subsec:limitations-linking-number}
%%%%%%%%%%%%%%%%%%%%%%%%%%%%%%%%%%%%%%%%%%%%%%%%%%%%%%%%%%%%

If

\[
\operatorname{lk}(L_1,L_2)\neq0,
\]

the components cannot be separated into an unlink.

However,

\[
\boxed{
\operatorname{lk}(L_1,L_2)=0
}
\]

does not imply that the link is trivial.

For example, the Whitehead link has linking number zero but is nontrivial.

Thus linking number is useful but incomplete.

%%%%%%%%%%%%%%%%%%%%%%%%%%%%%%%%%%%%%%%%%%%%%%%%%%%%%%%%%%%%
\subsection{Knot Complements}
\label{subsec:knot-complements}
%%%%%%%%%%%%%%%%%%%%%%%%%%%%%%%%%%%%%%%%%%%%%%%%%%%%%%%%%%%%

A knot can be studied through the space surrounding it.

Let

\[
K\subset S^3.
\]

One may consider

\[
S^3\setminus K.
\]

More commonly in \(3\)-manifold topology, one removes the interior of a
tubular neighborhood \(N(K)\):

\[
\boxed{
X_K
=
S^3\setminus\operatorname{int}N(K).
}
\]

This compact \(3\)-manifold is called the knot exterior.

Its boundary is

\[
\boxed{
\partial X_K\cong T^2.
}
\]

%%%%%%%%%%%%%%%%%%%%%%%%%%%%%%%%%%%%%%%%%%%%%%%%%%%%%%%%%%%%
\subsection{Why Knot Complements Matter}
\label{subsec:why-knot-complements}
%%%%%%%%%%%%%%%%%%%%%%%%%%%%%%%%%%%%%%%%%%%%%%%%%%%%%%%%%%%%

A knot is encoded remarkably strongly by its complement.

A foundational theorem states that knots in \(S^3\) are determined, up to
the appropriate equivalence, by the homeomorphism type of their complements.

Thus one can transform

\[
\boxed{
\text{knot problem}
}
\]

into

\[
\boxed{
\text{three-manifold problem}.
}
\]

%%%%%%%%%%%%%%%%%%%%%%%%%%%%%%%%%%%%%%%%%%%%%%%%%%%%%%%%%%%%
\subsection{The Knot Group}
\label{subsec:knot-group}
%%%%%%%%%%%%%%%%%%%%%%%%%%%%%%%%%%%%%%%%%%%%%%%%%%%%%%%%%%%%

\begin{definitionbox}{Knot Group}
The \emph{knot group} of \(K\) is

\[
\boxed{
G_K
=
\pi_1(S^3\setminus K).
}
\]
\end{definitionbox}

Equivalently, one may use the knot exterior:

\[
\boxed{
G_K
\cong
\pi_1(X_K).
}
\]

The knot group is a fundamental algebraic invariant.

%%%%%%%%%%%%%%%%%%%%%%%%%%%%%%%%%%%%%%%%%%%%%%%%%%%%%%%%%%%%
\subsection{Knot Group of the Unknot}
\label{subsec:unknot-group}
%%%%%%%%%%%%%%%%%%%%%%%%%%%%%%%%%%%%%%%%%%%%%%%%%%%%%%%%%%%%

The complement of the unknot has fundamental group

\[
\boxed{
\pi_1(S^3\setminus K_0)
\cong
\Z.
}
\]

Therefore the unknot group is infinite cyclic.

%%%%%%%%%%%%%%%%%%%%%%%%%%%%%%%%%%%%%%%%%%%%%%%%%%%%%%%%%%%%
\subsection{Wirtinger Presentation}
\label{subsec:wirtinger-presentation}
%%%%%%%%%%%%%%%%%%%%%%%%%%%%%%%%%%%%%%%%%%%%%%%%%%%%%%%%%%%%

A knot diagram provides a presentation of the knot group.

Assign a generator to each arc of the diagram.

At each crossing, impose a relation determined by the over-strand and
under-strands.

This construction is called the \emph{Wirtinger presentation}.

Schematically,

\[
\boxed{
G_K
=
\langle
x_1,\ldots,x_n
\mid
r_1,\ldots,r_n
\rangle.
}
\]

One relation is redundant for a connected knot diagram.

%%%%%%%%%%%%%%%%%%%%%%%%%%%%%%%%%%%%%%%%%%%%%%%%%%%%%%%%%%%%
\subsection{Trefoil Knot Group}
\label{subsec:trefoil-knot-group}
%%%%%%%%%%%%%%%%%%%%%%%%%%%%%%%%%%%%%%%%%%%%%%%%%%%%%%%%%%%%

The trefoil knot group admits the presentation

\[
\boxed{
G_{3_1}
=
\langle
a,b
\mid
aba=bab
\rangle.
}
\]

It also admits the presentation

\[
\boxed{
G_{3_1}
\cong
\langle
x,y
\mid
x^2=y^3
\rangle.
}
\]

This group is nonabelian.

Since

\[
\Z
\]

is abelian, the trefoil cannot be the unknot.

%%%%%%%%%%%%%%%%%%%%%%%%%%%%%%%%%%%%%%%%%%%%%%%%%%%%%%%%%%%%
\subsection{Worked Example: Distinguishing a Knot from the Unknot}
\label{subsec:worked-knot-group}
%%%%%%%%%%%%%%%%%%%%%%%%%%%%%%%%%%%%%%%%%%%%%%%%%%%%%%%%%%%%

\begin{workedexamplebox}{Using the Knot Group}

Suppose a knot \(K\) has knot group

\[
G_K
=
\langle
a,b
\mid
aba=bab
\rangle.
\]

This group is nonabelian.

The unknot group is

\[
G_{K_0}\cong\Z,
\]

which is abelian.

Therefore the groups cannot be isomorphic.

Hence \(K\) cannot be equivalent to the unknot.

\begin{answerbox}
\[
\boxed{
K\text{ is nontrivial}.
}
\]
\end{answerbox}
\end{workedexamplebox}

%%%%%%%%%%%%%%%%%%%%%%%%%%%%%%%%%%%%%%%%%%%%%%%%%%%%%%%%%%%%
\subsection{Abelianization of a Knot Group}
\label{subsec:abelianization-knot-group}
%%%%%%%%%%%%%%%%%%%%%%%%%%%%%%%%%%%%%%%%%%%%%%%%%%%%%%%%%%%%

For every knot \(K\subset S^3\),

\[
\boxed{
G_K^{\mathrm{ab}}
\cong
\Z.
}
\]

Equivalently,

\[
\boxed{
H_1(S^3\setminus K;\Z)
\cong
\Z.
}
\]

Thus ordinary first homology cannot distinguish knots from one another.

More refined information must be extracted from the nonabelian group or from
additional constructions.

%%%%%%%%%%%%%%%%%%%%%%%%%%%%%%%%%%%%%%%%%%%%%%%%%%%%%%%%%%%%
\subsection{Knot Colorings}
\label{subsec:knot-colorings}
%%%%%%%%%%%%%%%%%%%%%%%%%%%%%%%%%%%%%%%%%%%%%%%%%%%%%%%%%%%%

Certain diagram colorings provide elementary knot invariants.

A classical example is Fox \(n\)-coloring.

Assign elements of

\[
\Z/n\Z
\]

to arcs of a knot diagram.

At each crossing, require

\[
\boxed{
2b\equiv a+c\pmod n,
}
\]

where \(b\) is the color of the over-arc and \(a,c\) are the colors of the
under-arcs.

%%%%%%%%%%%%%%%%%%%%%%%%%%%%%%%%%%%%%%%%%%%%%%%%%%%%%%%%%%%%
\subsection{Fox Three-Coloring}
\label{subsec:fox-three-coloring}
%%%%%%%%%%%%%%%%%%%%%%%%%%%%%%%%%%%%%%%%%%%%%%%%%%%%%%%%%%%%

For

\[
n=3,
\]

the rule implies that at every crossing either

\[
\boxed{
\text{all three colors agree}
}
\]

or

\[
\boxed{
\text{all three colors are different}.
}
\]

A nontrivial coloring uses more than one color.

The existence of a nontrivial \(3\)-coloring is preserved by Reidemeister
moves.

%%%%%%%%%%%%%%%%%%%%%%%%%%%%%%%%%%%%%%%%%%%%%%%%%%%%%%%%%%%%
\subsection{Trefoil Three-Colorability}
\label{subsec:trefoil-three-coloring}
%%%%%%%%%%%%%%%%%%%%%%%%%%%%%%%%%%%%%%%%%%%%%%%%%%%%%%%%%%%%

The trefoil admits a nontrivial \(3\)-coloring.

The unknot does not admit a nontrivial \(3\)-coloring.

Therefore three-colorability proves

\[
\boxed{
\text{trefoil}\not\cong\text{unknot}.
}
\]

This gives an accessible example of a combinatorial knot invariant.

%%%%%%%%%%%%%%%%%%%%%%%%%%%%%%%%%%%%%%%%%%%%%%%%%%%%%%%%%%%%
\subsection{Polynomial Invariants}
\label{subsec:polynomial-knot-invariants}
%%%%%%%%%%%%%%%%%%%%%%%%%%%%%%%%%%%%%%%%%%%%%%%%%%%%%%%%%%%%

One of the major ideas in knot theory is to assign a polynomial to each knot.

If equivalent knots have the same polynomial, then the polynomial is an
invariant.

Important examples include

\[
\boxed{
\text{Alexander polynomial},
}
\]

\[
\boxed{
\text{Jones polynomial},
}
\]

and

\[
\boxed{
\text{HOMFLY-PT polynomial}.
}
\]

%%%%%%%%%%%%%%%%%%%%%%%%%%%%%%%%%%%%%%%%%%%%%%%%%%%%%%%%%%%%
\subsection{Alexander Polynomial}
\label{subsec:alexander-polynomial}
%%%%%%%%%%%%%%%%%%%%%%%%%%%%%%%%%%%%%%%%%%%%%%%%%%%%%%%%%%%%

The Alexander polynomial of a knot \(K\) is denoted

\[
\boxed{
\Delta_K(t).
}
\]

The uppercase Greek letter

\[
\Delta
\]

is read ``Delta.''

The polynomial is defined only up to multiplication by a unit

\[
\boxed{
\pm t^m.
}
\]

A common normalization imposes additional conditions to choose a preferred
representative.

%%%%%%%%%%%%%%%%%%%%%%%%%%%%%%%%%%%%%%%%%%%%%%%%%%%%%%%%%%%%
\subsection{Examples of Alexander Polynomials}
\label{subsec:alexander-polynomial-examples}
%%%%%%%%%%%%%%%%%%%%%%%%%%%%%%%%%%%%%%%%%%%%%%%%%%%%%%%%%%%%

For the unknot,

\[
\boxed{
\Delta_{K_0}(t)=1.
}
\]

For the trefoil, one common normalized form is

\[
\boxed{
\Delta_{3_1}(t)
=
t^2-t+1.
}
\]

Equivalently, up to multiplication by a unit,

\[
t-1+t^{-1}
\]

represents the same Alexander polynomial.

%%%%%%%%%%%%%%%%%%%%%%%%%%%%%%%%%%%%%%%%%%%%%%%%%%%%%%%%%%%%
\subsection{Alexander Polynomial and Connected Sum}
\label{subsec:alexander-connected-sum}
%%%%%%%%%%%%%%%%%%%%%%%%%%%%%%%%%%%%%%%%%%%%%%%%%%%%%%%%%%%%

The Alexander polynomial satisfies

\[
\boxed{
\Delta_{K_1\#K_2}(t)
=
\Delta_{K_1}(t)\Delta_{K_2}(t)
}
\]

up to the usual unit ambiguity.

Thus connected sum becomes multiplication at the polynomial level.

%%%%%%%%%%%%%%%%%%%%%%%%%%%%%%%%%%%%%%%%%%%%%%%%%%%%%%%%%%%%
\subsection{Limitations of the Alexander Polynomial}
\label{subsec:limitations-alexander-polynomial}
%%%%%%%%%%%%%%%%%%%%%%%%%%%%%%%%%%%%%%%%%%%%%%%%%%%%%%%%%%%%

The Alexander polynomial does not completely classify knots.

Different knots can have the same Alexander polynomial.

In particular,

\[
\boxed{
\Delta_K(t)=1
}
\]

does not imply that \(K\) is the unknot.

This illustrates an important principle:

\[
\boxed{
\text{an invariant can distinguish some objects without classifying all objects}.
}
\]

%%%%%%%%%%%%%%%%%%%%%%%%%%%%%%%%%%%%%%%%%%%%%%%%%%%%%%%%%%%%
\subsection{Jones Polynomial}
\label{subsec:jones-polynomial}
%%%%%%%%%%%%%%%%%%%%%%%%%%%%%%%%%%%%%%%%%%%%%%%%%%%%%%%%%%%%

The Jones polynomial is another polynomial invariant.

It is commonly denoted

\[
\boxed{
V_K(t).
}
\]

It emerged from connections between knot theory, statistical mechanics,
operator algebras, and braid representations.

The Jones polynomial is particularly sensitive to crossing information.

%%%%%%%%%%%%%%%%%%%%%%%%%%%%%%%%%%%%%%%%%%%%%%%%%%%%%%%%%%%%
\subsection{Skein Relations}
\label{subsec:skein-relations}
%%%%%%%%%%%%%%%%%%%%%%%%%%%%%%%%%%%%%%%%%%%%%%%%%%%%%%%%%%%%

Many polynomial invariants can be defined recursively through local changes
at a crossing.

One considers three diagrams

\[
L_+,
\qquad
L_-,
\qquad
L_0,
\]

which differ only inside a small disk.

The diagrams \(L_+\) and \(L_-\) contain positive and negative crossings,
while \(L_0\) contains a smoothing.

A relation among their polynomials is called a \emph{skein relation}.

%%%%%%%%%%%%%%%%%%%%%%%%%%%%%%%%%%%%%%%%%%%%%%%%%%%%%%%%%%%%
\subsection{Alexander Skein Relation}
\label{subsec:alexander-skein-relation}
%%%%%%%%%%%%%%%%%%%%%%%%%%%%%%%%%%%%%%%%%%%%%%%%%%%%%%%%%%%%

Under a common normalization, the Alexander polynomial satisfies

\[
\boxed{
\Delta_{L_+}(t)
-
\Delta_{L_-}(t)
=
\left(
t^{1/2}-t^{-1/2}
\right)
\Delta_{L_0}(t).
}
\]

Different conventions may use slightly different normalizations.

The essential idea is that a local crossing change relates the invariants of
three links.

%%%%%%%%%%%%%%%%%%%%%%%%%%%%%%%%%%%%%%%%%%%%%%%%%%%%%%%%%%%%
\subsection{Jones Skein Relation}
\label{subsec:jones-skein-relation}
%%%%%%%%%%%%%%%%%%%%%%%%%%%%%%%%%%%%%%%%%%%%%%%%%%%%%%%%%%%%

One common normalization of the Jones polynomial satisfies

\[
\boxed{
t^{-1}V_{L_+}(t)
-
tV_{L_-}(t)
=
\left(
t^{1/2}-t^{-1/2}
\right)V_{L_0}(t).
}
\]

Together with

\[
\boxed{
V_{\text{unknot}}(t)=1,
}
\]

the skein relation recursively determines the polynomial.

%%%%%%%%%%%%%%%%%%%%%%%%%%%%%%%%%%%%%%%%%%%%%%%%%%%%%%%%%%%%
\subsection{HOMFLY-PT Polynomial}
\label{subsec:homfly-pt}
%%%%%%%%%%%%%%%%%%%%%%%%%%%%%%%%%%%%%%%%%%%%%%%%%%%%%%%%%%%%

The HOMFLY-PT polynomial is a two-variable link invariant that generalizes
features of both the Alexander and Jones polynomials.

It is defined through a skein relation of the general form

\[
\boxed{
aP(L_+)-a^{-1}P(L_-)
=
zP(L_0),
}
\]

under a standard choice of variables and normalization.

Specializations recover other important polynomial invariants.

%%%%%%%%%%%%%%%%%%%%%%%%%%%%%%%%%%%%%%%%%%%%%%%%%%%%%%%%%%%%
\subsection{Seifert Surfaces}
\label{subsec:seifert-surfaces}
%%%%%%%%%%%%%%%%%%%%%%%%%%%%%%%%%%%%%%%%%%%%%%%%%%%%%%%%%%%%

\begin{definitionbox}{Seifert Surface}
A \emph{Seifert surface} for an oriented knot \(K\) is a compact connected
orientable surface \(F\subset S^3\) whose boundary is

\[
\boxed{
\partial F=K.
}
\]
\end{definitionbox}

Thus a knot can be studied through a surface that it bounds.

%%%%%%%%%%%%%%%%%%%%%%%%%%%%%%%%%%%%%%%%%%%%%%%%%%%%%%%%%%%%
\subsection{Existence of Seifert Surfaces}
\label{subsec:existence-seifert-surfaces}
%%%%%%%%%%%%%%%%%%%%%%%%%%%%%%%%%%%%%%%%%%%%%%%%%%%%%%%%%%%%

Every oriented knot in \(S^3\) bounds a Seifert surface.

Therefore

\[
\boxed{
\text{knot}
\longrightarrow
\text{bounding surface}
}
\]

provides another bridge from one-dimensional embeddings to
two-dimensional topology.

%%%%%%%%%%%%%%%%%%%%%%%%%%%%%%%%%%%%%%%%%%%%%%%%%%%%%%%%%%%%
\subsection{Knot Genus}
\label{subsec:knot-genus}
%%%%%%%%%%%%%%%%%%%%%%%%%%%%%%%%%%%%%%%%%%%%%%%%%%%%%%%%%%%%

\begin{definitionbox}{Knot Genus}
The \emph{genus} of a knot \(K\), denoted

\[
\boxed{
g(K),
}
\]

is the minimum genus among all Seifert surfaces for \(K\).
\end{definitionbox}

For the unknot,

\[
\boxed{
g(K_0)=0.
}
\]

For the trefoil,

\[
\boxed{
g(3_1)=1.
}
\]

%%%%%%%%%%%%%%%%%%%%%%%%%%%%%%%%%%%%%%%%%%%%%%%%%%%%%%%%%%%%
\subsection{Genus Detects the Unknot}
\label{subsec:genus-detects-unknot}
%%%%%%%%%%%%%%%%%%%%%%%%%%%%%%%%%%%%%%%%%%%%%%%%%%%%%%%%%%%%

A knot has genus zero if and only if it is the unknot:

\[
\boxed{
g(K)=0
\Longleftrightarrow
K\text{ is the unknot}.
}
\]

Indeed, a connected orientable genus-zero surface with one boundary
component is a disk.

Thus a genus-zero knot bounds an embedded disk in \(S^3\).

%%%%%%%%%%%%%%%%%%%%%%%%%%%%%%%%%%%%%%%%%%%%%%%%%%%%%%%%%%%%
\subsection{Seifert Algorithm}
\label{subsec:seifert-algorithm}
%%%%%%%%%%%%%%%%%%%%%%%%%%%%%%%%%%%%%%%%%%%%%%%%%%%%%%%%%%%%

A Seifert surface can be constructed directly from an oriented knot diagram.

The basic procedure is:

\begin{enumerate}
    \item orient the knot;
    \item smooth every crossing according to the orientation;
    \item obtain a collection of Seifert circles;
    \item fill each circle with a disk;
    \item reconnect the disks using twisted bands corresponding to the
    original crossings.
\end{enumerate}

The result is an orientable surface bounded by the knot.

%%%%%%%%%%%%%%%%%%%%%%%%%%%%%%%%%%%%%%%%%%%%%%%%%%%%%%%%%%%%
\subsection{Genus from the Seifert Algorithm}
\label{subsec:genus-seifert-algorithm}
%%%%%%%%%%%%%%%%%%%%%%%%%%%%%%%%%%%%%%%%%%%%%%%%%%%%%%%%%%%%

Suppose a knot diagram has

\[
c
\]

crossings and produces

\[
s
\]

Seifert circles.

The resulting Seifert surface has Euler characteristic

\[
\boxed{
\chi=s-c.
}
\]

Since the surface has one boundary component,

\[
\chi=1-2g.
\]

Therefore the genus of the constructed surface is

\[
\boxed{
g
=
\frac{1-s+c}{2}.
}
\]

This gives an upper bound for the knot genus.

%%%%%%%%%%%%%%%%%%%%%%%%%%%%%%%%%%%%%%%%%%%%%%%%%%%%%%%%%%%%
\subsection{Worked Example: Seifert Surface Genus}
\label{subsec:worked-seifert-genus}
%%%%%%%%%%%%%%%%%%%%%%%%%%%%%%%%%%%%%%%%%%%%%%%%%%%%%%%%%%%%

\begin{workedexamplebox}{Genus from a Diagram}

Suppose an oriented knot diagram has

\[
c=7
\]

crossings and

\[
s=4
\]

Seifert circles.

Then the Seifert surface produced by the algorithm has

\[
\chi=s-c=4-7=-3.
\]

For a connected orientable surface with one boundary component,

\[
\chi=1-2g.
\]

Hence

\[
-3=1-2g.
\]

Therefore

\[
g=2.
\]

\begin{answerbox}
\[
\boxed{
g(F)=2.
}
\]

Thus

\[
\boxed{
g(K)\leq2.
}
\]
\end{answerbox}
\end{workedexamplebox}

%%%%%%%%%%%%%%%%%%%%%%%%%%%%%%%%%%%%%%%%%%%%%%%%%%%%%%%%%%%%
\subsection{Seifert Matrices}
\label{subsec:seifert-matrices}
%%%%%%%%%%%%%%%%%%%%%%%%%%%%%%%%%%%%%%%%%%%%%%%%%%%%%%%%%%%%

Choose a basis

\[
a_1,\ldots,a_{2g}
\]

for the first homology of a Seifert surface.

Push one curve slightly off the surface in the positive normal direction.

The Seifert matrix \(V\) is defined by linking numbers:

\[
\boxed{
V_{ij}
=
\operatorname{lk}(a_i^+,a_j).
}
\]

This matrix contains substantial information about the knot.

%%%%%%%%%%%%%%%%%%%%%%%%%%%%%%%%%%%%%%%%%%%%%%%%%%%%%%%%%%%%
\subsection{Alexander Polynomial from a Seifert Matrix}
\label{subsec:alexander-seifert-matrix}
%%%%%%%%%%%%%%%%%%%%%%%%%%%%%%%%%%%%%%%%%%%%%%%%%%%%%%%%%%%%

The Alexander polynomial can be obtained from a Seifert matrix by

\[
\boxed{
\Delta_K(t)
\doteq
\det(V-tV^{T}),
}
\]

where

\[
\doteq
\]

means equality up to multiplication by a unit

\[
\pm t^m.
\]

This connects topology of surfaces, linking numbers, linear algebra, and
polynomial invariants.

%%%%%%%%%%%%%%%%%%%%%%%%%%%%%%%%%%%%%%%%%%%%%%%%%%%%%%%%%%%%
\subsection{Braids}
\label{subsec:braids-knot-theory}
%%%%%%%%%%%%%%%%%%%%%%%%%%%%%%%%%%%%%%%%%%%%%%%%%%%%%%%%%%%%

A braid consists of strands moving monotonically from one collection of
points to another.

Braids can be multiplied by stacking one braid on top of another.

The resulting algebraic structure is the braid group.

%%%%%%%%%%%%%%%%%%%%%%%%%%%%%%%%%%%%%%%%%%%%%%%%%%%%%%%%%%%%
\subsection{Braid Group}
\label{subsec:braid-group}
%%%%%%%%%%%%%%%%%%%%%%%%%%%%%%%%%%%%%%%%%%%%%%%%%%%%%%%%%%%%

The braid group on \(n\) strands is denoted

\[
\boxed{
B_n.
}
\]

It has generators

\[
\boxed{
\sigma_1,\ldots,\sigma_{n-1}.
}
\]

The lowercase Greek letter

\[
\sigma
\]

is read ``sigma.''

The generator \(\sigma_i\) represents the \(i\)-th strand crossing over the
\((i+1)\)-st strand.

%%%%%%%%%%%%%%%%%%%%%%%%%%%%%%%%%%%%%%%%%%%%%%%%%%%%%%%%%%%%
\subsection{Braid Relations}
\label{subsec:braid-relations}
%%%%%%%%%%%%%%%%%%%%%%%%%%%%%%%%%%%%%%%%%%%%%%%%%%%%%%%%%%%%

The braid group satisfies

\[
\boxed{
\sigma_i\sigma_j
=
\sigma_j\sigma_i
\qquad
|i-j|\geq2,
}
\]

and

\[
\boxed{
\sigma_i\sigma_{i+1}\sigma_i
=
\sigma_{i+1}\sigma_i\sigma_{i+1}.
}
\]

Thus

\[
\boxed{
B_n
=
\left\langle
\sigma_1,\ldots,\sigma_{n-1}
\mid
\text{braid relations}
\right\rangle.
}
\]

%%%%%%%%%%%%%%%%%%%%%%%%%%%%%%%%%%%%%%%%%%%%%%%%%%%%%%%%%%%%
\subsection{Braid Groups Are Nonabelian}
\label{subsec:braid-groups-nonabelian}
%%%%%%%%%%%%%%%%%%%%%%%%%%%%%%%%%%%%%%%%%%%%%%%%%%%%%%%%%%%%

For

\[
n\geq3,
\]

the braid group is nonabelian.

For example, in \(B_3\),

\[
\sigma_1\sigma_2
\neq
\sigma_2\sigma_1.
\]

However,

\[
\boxed{
\sigma_1\sigma_2\sigma_1
=
\sigma_2\sigma_1\sigma_2.
}
\]

This is the fundamental braid relation.

%%%%%%%%%%%%%%%%%%%%%%%%%%%%%%%%%%%%%%%%%%%%%%%%%%%%%%%%%%%%
\subsection{Closing a Braid}
\label{subsec:braid-closure}
%%%%%%%%%%%%%%%%%%%%%%%%%%%%%%%%%%%%%%%%%%%%%%%%%%%%%%%%%%%%

A braid can be converted into a link by connecting corresponding upper and
lower endpoints.

The resulting link is called the closure of the braid.

Thus

\[
\boxed{
\text{braid}
\longrightarrow
\text{link}.
}
\]

%%%%%%%%%%%%%%%%%%%%%%%%%%%%%%%%%%%%%%%%%%%%%%%%%%%%%%%%%%%%
\subsection{Alexander's Theorem for Braids}
\label{subsec:alexander-theorem-braids}
%%%%%%%%%%%%%%%%%%%%%%%%%%%%%%%%%%%%%%%%%%%%%%%%%%%%%%%%%%%%

\begin{theorem}[Alexander's Theorem]
Every oriented link in \(S^3\) is isotopic to the closure of some braid.
\end{theorem}

Thus braid theory provides an algebraic framework capable of representing
every link.

%%%%%%%%%%%%%%%%%%%%%%%%%%%%%%%%%%%%%%%%%%%%%%%%%%%%%%%%%%%%
\subsection{Markov Moves}
\label{subsec:markov-moves}
%%%%%%%%%%%%%%%%%%%%%%%%%%%%%%%%%%%%%%%%%%%%%%%%%%%%%%%%%%%%

Different braids can have isotopic closures.

Markov's theorem describes the operations needed to relate braid
representatives of the same oriented link.

Conceptually, the two basic operations are

\[
\boxed{
\text{conjugation}
}
\]

and

\[
\boxed{
\text{stabilization/destabilization}.
}
\]

Together with braid relations, these moves translate link equivalence into
an algebraic problem involving braid groups.

%%%%%%%%%%%%%%%%%%%%%%%%%%%%%%%%%%%%%%%%%%%%%%%%%%%%%%%%%%%%
\subsection{Torus Knots}
\label{subsec:torus-knots}
%%%%%%%%%%%%%%%%%%%%%%%%%%%%%%%%%%%%%%%%%%%%%%%%%%%%%%%%%%%%

A torus knot lies on the surface of an unknotted torus.

It is denoted

\[
\boxed{
T(p,q),
}
\]

where the curve winds \(p\) times in one toroidal direction and \(q\) times
in the other.

When

\[
\gcd(p,q)=1,
\]

the result is a knot.

If

\[
\gcd(p,q)>1,
\]

the result is a link.

%%%%%%%%%%%%%%%%%%%%%%%%%%%%%%%%%%%%%%%%%%%%%%%%%%%%%%%%%%%%
\subsection{Trefoil as a Torus Knot}
\label{subsec:trefoil-torus-knot}
%%%%%%%%%%%%%%%%%%%%%%%%%%%%%%%%%%%%%%%%%%%%%%%%%%%%%%%%%%%%

The trefoil can be represented as

\[
\boxed{
T(2,3).
}
\]

Up to orientation and convention, one may also encounter

\[
T(3,2).
\]

This makes the trefoil the simplest nontrivial torus knot.

%%%%%%%%%%%%%%%%%%%%%%%%%%%%%%%%%%%%%%%%%%%%%%%%%%%%%%%%%%%%
\subsection{Crossing Number of Torus Knots}
\label{subsec:crossing-number-torus-knots}
%%%%%%%%%%%%%%%%%%%%%%%%%%%%%%%%%%%%%%%%%%%%%%%%%%%%%%%%%%%%

For positive coprime integers \(p\) and \(q\),

\[
\boxed{
c(T(p,q))
=
\min\{(p-1)q,(q-1)p\}.
}
\]

For the trefoil,

\[
T(2,3),
\]

we obtain

\[
c(T(2,3))
=
\min\{3,4\}
=
3.
\]

%%%%%%%%%%%%%%%%%%%%%%%%%%%%%%%%%%%%%%%%%%%%%%%%%%%%%%%%%%%%
\subsection{Genus of a Torus Knot}
\label{subsec:genus-torus-knot}
%%%%%%%%%%%%%%%%%%%%%%%%%%%%%%%%%%%%%%%%%%%%%%%%%%%%%%%%%%%%

For coprime positive integers \(p,q\),

\[
\boxed{
g(T(p,q))
=
\frac{(p-1)(q-1)}{2}.
}
\]

For the trefoil,

\[
g(T(2,3))
=
\frac{(2-1)(3-1)}{2}
=
1.
\]

%%%%%%%%%%%%%%%%%%%%%%%%%%%%%%%%%%%%%%%%%%%%%%%%%%%%%%%%%%%%
\subsection{Alternating Knots}
\label{subsec:alternating-knots}
%%%%%%%%%%%%%%%%%%%%%%%%%%%%%%%%%%%%%%%%%%%%%%%%%%%%%%%%%%%%

A knot diagram is alternating if, while traveling along the diagram, the
crossings alternate between

\[
\boxed{
\text{over}
\quad\text{and}\quad
\text{under}.
}
\]

A knot is alternating if it possesses an alternating diagram.

The trefoil and figure-eight knot are alternating.

Alternating knots possess especially strong relationships between their
diagrams and several knot invariants.

%%%%%%%%%%%%%%%%%%%%%%%%%%%%%%%%%%%%%%%%%%%%%%%%%%%%%%%%%%%%
\subsection{Reduced Alternating Diagrams}
\label{subsec:reduced-alternating-diagrams}
%%%%%%%%%%%%%%%%%%%%%%%%%%%%%%%%%%%%%%%%%%%%%%%%%%%%%%%%%%%%

A reduced alternating diagram contains no removable nugatory crossing.

A fundamental result in knot theory implies that a reduced alternating
diagram realizes the crossing number of the represented link.

Thus for this important class of diagrams,

\[
\boxed{
\text{visible crossings}
=
\text{minimum crossings}.
}
\]

%%%%%%%%%%%%%%%%%%%%%%%%%%%%%%%%%%%%%%%%%%%%%%%%%%%%%%%%%%%%
\subsection{Knot Determinant}
\label{subsec:knot-determinant}
%%%%%%%%%%%%%%%%%%%%%%%%%%%%%%%%%%%%%%%%%%%%%%%%%%%%%%%%%%%%

The determinant of a knot is defined by

\[
\boxed{
\det(K)
=
|\Delta_K(-1)|.
}
\]

For the trefoil,

\[
\Delta_{3_1}(t)
=
t^2-t+1.
\]

Therefore

\[
\Delta_{3_1}(-1)
=
1+1+1
=
3,
\]

so

\[
\boxed{
\det(3_1)=3.
}
\]

%%%%%%%%%%%%%%%%%%%%%%%%%%%%%%%%%%%%%%%%%%%%%%%%%%%%%%%%%%%%
\subsection{Determinant and Colorings}
\label{subsec:determinant-colorings}
%%%%%%%%%%%%%%%%%%%%%%%%%%%%%%%%%%%%%%%%%%%%%%%%%%%%%%%%%%%%

The knot determinant is related to Fox colorings.

In particular, nontrivial Fox \(p\)-colorings for a prime \(p\) are closely
controlled by whether

\[
\boxed{
p\mid\det(K).
}
\]

For the trefoil,

\[
\det(3_1)=3,
\]

which is consistent with its nontrivial three-colorability.

%%%%%%%%%%%%%%%%%%%%%%%%%%%%%%%%%%%%%%%%%%%%%%%%%%%%%%%%%%%%
\subsection{Writhe}
\label{subsec:knot-writhe}
%%%%%%%%%%%%%%%%%%%%%%%%%%%%%%%%%%%%%%%%%%%%%%%%%%%%%%%%%%%%

For an oriented knot diagram \(D\), the writhe is

\[
\boxed{
w(D)
=
\sum_c\operatorname{sgn}(c).
}
\]

Writhe is preserved by Reidemeister moves II and III but changes under
Reidemeister move I.

Therefore writhe alone is not a knot invariant.

It is a diagram invariant under regular isotopy.

%%%%%%%%%%%%%%%%%%%%%%%%%%%%%%%%%%%%%%%%%%%%%%%%%%%%%%%%%%%%
\subsection{Kauffman Bracket}
\label{subsec:kauffman-bracket}
%%%%%%%%%%%%%%%%%%%%%%%%%%%%%%%%%%%%%%%%%%%%%%%%%%%%%%%%%%%%

The Kauffman bracket assigns a Laurent polynomial

\[
\boxed{
\langle D\rangle
}
\]

to an unoriented link diagram.

It is defined recursively through smoothing rules at crossings.

The bracket is invariant under Reidemeister moves II and III but not move I.

After a normalization involving writhe, it yields the Jones polynomial.

%%%%%%%%%%%%%%%%%%%%%%%%%%%%%%%%%%%%%%%%%%%%%%%%%%%%%%%%%%%%
\subsection{From Diagram Invariants to Knot Invariants}
\label{subsec:diagram-to-knot-invariant}
%%%%%%%%%%%%%%%%%%%%%%%%%%%%%%%%%%%%%%%%%%%%%%%%%%%%%%%%%%%%

A common strategy in knot theory is:

\[
\boxed{
\text{define quantity on diagrams}
}
\]

then prove invariance under

\[
\boxed{
R_1,\quad R_2,\quad R_3.
}
\]

Once invariance under all Reidemeister moves is established, the quantity
descends to an invariant of knots.

This principle explains why Reidemeister moves are so central.

%%%%%%%%%%%%%%%%%%%%%%%%%%%%%%%%%%%%%%%%%%%%%%%%%%%%%%%%%%%%
\subsection{Hyperbolic Knots}
\label{subsec:hyperbolic-knots}
%%%%%%%%%%%%%%%%%%%%%%%%%%%%%%%%%%%%%%%%%%%%%%%%%%%%%%%%%%%%

A knot is called hyperbolic if the interior of its knot exterior admits a
complete finite-volume hyperbolic metric.

Thus

\[
\boxed{
S^3\setminus K
}
\]

carries hyperbolic geometry.

The figure-eight knot is the standard elementary example of a hyperbolic
knot.

%%%%%%%%%%%%%%%%%%%%%%%%%%%%%%%%%%%%%%%%%%%%%%%%%%%%%%%%%%%%
\subsection{Hyperbolic Volume}
\label{subsec:knot-hyperbolic-volume}
%%%%%%%%%%%%%%%%%%%%%%%%%%%%%%%%%%%%%%%%%%%%%%%%%%%%%%%%%%%%

For a hyperbolic knot \(K\), the hyperbolic volume

\[
\boxed{
\operatorname{Vol}(S^3\setminus K)
}
\]

is a knot invariant.

Mostow--Prasad rigidity ensures that the finite-volume hyperbolic structure
is rigid.

Thus a geometric quantity becomes a topological invariant.

%%%%%%%%%%%%%%%%%%%%%%%%%%%%%%%%%%%%%%%%%%%%%%%%%%%%%%%%%%%%
\subsection{Torus, Satellite, and Hyperbolic Knots}
\label{subsec:torus-satellite-hyperbolic}
%%%%%%%%%%%%%%%%%%%%%%%%%%%%%%%%%%%%%%%%%%%%%%%%%%%%%%%%%%%%

A major structural viewpoint divides prime knots into three broad geometric
types:

\[
\boxed{
\text{torus knots},
}
\]

\[
\boxed{
\text{satellite knots},
}
\]

and

\[
\boxed{
\text{hyperbolic knots}.
}
\]

This classification is closely related to the decomposition theory of knot
complements and the geometrization of \(3\)-manifolds.

%%%%%%%%%%%%%%%%%%%%%%%%%%%%%%%%%%%%%%%%%%%%%%%%%%%%%%%%%%%%
\subsection{Satellite Knots}
\label{subsec:satellite-knots}
%%%%%%%%%%%%%%%%%%%%%%%%%%%%%%%%%%%%%%%%%%%%%%%%%%%%%%%%%%%%

A satellite knot is constructed by placing a nontrivial knot pattern inside
a solid torus and then embedding that solid torus as a tubular neighborhood
of another knot.

Satellite knots therefore contain nested geometric structure.

Their complements contain essential tori, distinguishing them from
hyperbolic knots.

%%%%%%%%%%%%%%%%%%%%%%%%%%%%%%%%%%%%%%%%%%%%%%%%%%%%%%%%%%%%
\subsection{Dehn Surgery}
\label{subsec:knot-dehn-surgery}
%%%%%%%%%%%%%%%%%%%%%%%%%%%%%%%%%%%%%%%%%%%%%%%%%%%%%%%%%%%%

A knot exterior has torus boundary.

Choose a slope on

\[
\partial X_K\cong T^2
\]

and glue in a solid torus so that the chosen slope becomes the meridian.

The resulting space is a new \(3\)-manifold.

This operation is called Dehn surgery or Dehn filling.

Schematically,

\[
\boxed{
\text{knot}
\longrightarrow
\text{knot exterior}
\longrightarrow
\text{new }3\text{-manifold}.
}
\]

%%%%%%%%%%%%%%%%%%%%%%%%%%%%%%%%%%%%%%%%%%%%%%%%%%%%%%%%%%%%
\subsection{Meridians and Longitudes}
\label{subsec:knot-meridian-longitude}
%%%%%%%%%%%%%%%%%%%%%%%%%%%%%%%%%%%%%%%%%%%%%%%%%%%%%%%%%%%%

The boundary torus of a knot exterior contains two important curve classes.

The \emph{meridian}

\[
\boxed{
\mu
}
\]

bounds a disk in the removed tubular neighborhood.

The \emph{longitude}

\[
\boxed{
\lambda
}
\]

runs parallel to the knot and is chosen with zero linking number with the
knot under the standard preferred-longitude convention.

Together,

\[
\boxed{
(\mu,\lambda)
}
\]

form a basis for

\[
H_1(T^2;\Z)\cong\Z^2.
\]

%%%%%%%%%%%%%%%%%%%%%%%%%%%%%%%%%%%%%%%%%%%%%%%%%%%%%%%%%%%%
\subsection{Surgery Slopes}
\label{subsec:knot-surgery-slopes}
%%%%%%%%%%%%%%%%%%%%%%%%%%%%%%%%%%%%%%%%%%%%%%%%%%%%%%%%%%%%

A primitive slope on the boundary torus has the form

\[
\boxed{
p\mu+q\lambda,
}
\]

where

\[
\gcd(p,q)=1.
\]

Different slopes generally produce different \(3\)-manifolds.

Thus knots provide a systematic mechanism for constructing manifolds.

%%%%%%%%%%%%%%%%%%%%%%%%%%%%%%%%%%%%%%%%%%%%%%%%%%%%%%%%%%%%
\subsection{Concordance}
\label{subsec:knot-concordance}
%%%%%%%%%%%%%%%%%%%%%%%%%%%%%%%%%%%%%%%%%%%%%%%%%%%%%%%%%%%%

Knot concordance weakens ordinary knot equivalence by allowing knots to be
connected through an annulus in four dimensions.

Two knots \(K_0\) and \(K_1\) are smoothly concordant if there exists a
smooth embedding

\[
\boxed{
S^1\times[0,1]
\hookrightarrow
S^3\times[0,1]
}
\]

whose boundary is the two knots at the ends.

Thus concordance introduces a fundamentally four-dimensional viewpoint.

%%%%%%%%%%%%%%%%%%%%%%%%%%%%%%%%%%%%%%%%%%%%%%%%%%%%%%%%%%%%
\subsection{Slice Knots}
\label{subsec:slice-knots}
%%%%%%%%%%%%%%%%%%%%%%%%%%%%%%%%%%%%%%%%%%%%%%%%%%%%%%%%%%%%

A knot \(K\subset S^3\) is slice if it bounds a smoothly embedded disk in the
\(4\)-ball:

\[
\boxed{
D^2\hookrightarrow D^4,
\qquad
\partial D^2=K.
}
\]

A knot may fail to bound a disk in \(S^3\) but succeed in \(D^4\).

Therefore knot theory changes substantially when one moves from three to
four dimensions.

%%%%%%%%%%%%%%%%%%%%%%%%%%%%%%%%%%%%%%%%%%%%%%%%%%%%%%%%%%%%
\subsection{Concordance Group}
\label{subsec:knot-concordance-group}
%%%%%%%%%%%%%%%%%%%%%%%%%%%%%%%%%%%%%%%%%%%%%%%%%%%%%%%%%%%%

Smooth concordance classes of oriented knots form an abelian group under
connected sum.

The identity is represented by the unknot.

The inverse of a knot is obtained using the appropriate mirror image and
orientation reversal.

This group is called the smooth knot concordance group.

%%%%%%%%%%%%%%%%%%%%%%%%%%%%%%%%%%%%%%%%%%%%%%%%%%%%%%%%%%%%
\subsection{Knot Invariants and Four Dimensions}
\label{subsec:knot-invariants-four-dimensions}
%%%%%%%%%%%%%%%%%%%%%%%%%%%%%%%%%%%%%%%%%%%%%%%%%%%%%%%%%%%%

Many invariants can obstruct a knot from being slice.

Classical examples arise from

\[
\boxed{
\text{Seifert forms},
\quad
\text{signatures},
\quad
\text{Alexander polynomial}.
}
\]

Modern invariants arise from deeper theories such as

\[
\boxed{
\text{gauge theory},
\quad
\text{Floer homology},
\quad
\text{Khovanov homology}.
}
\]

These theories illustrate the continuing interaction between knot theory,
topology, algebra, and geometry.

%%%%%%%%%%%%%%%%%%%%%%%%%%%%%%%%%%%%%%%%%%%%%%%%%%%%%%%%%%%%
\subsection{Spatial Graphs}
\label{subsec:spatial-graphs}
%%%%%%%%%%%%%%%%%%%%%%%%%%%%%%%%%%%%%%%%%%%%%%%%%%%%%%%%%%%%

Knot theory can be generalized by embedding graphs in three-dimensional
space.

A spatial graph is an embedding

\[
\boxed{
G\hookrightarrow S^3
}
\]

of a graph \(G\).

Knots correspond to the special case

\[
G=S^1.
\]

Spatial graph theory studies linking and knotting phenomena involving
vertices and edges.

%%%%%%%%%%%%%%%%%%%%%%%%%%%%%%%%%%%%%%%%%%%%%%%%%%%%%%%%%%%%
\subsection{Higher-Dimensional Knots}
\label{subsec:higher-dimensional-knots}
%%%%%%%%%%%%%%%%%%%%%%%%%%%%%%%%%%%%%%%%%%%%%%%%%%%%%%%%%%%%

The concept of knotting extends beyond circles in three-space.

A higher-dimensional knot may be represented by an embedding

\[
\boxed{
S^n\hookrightarrow S^{n+2}.
}
\]

Classical knot theory corresponds to

\[
n=1.
\]

Thus

\[
S^1\hookrightarrow S^3
\]

is the first case of a much broader theory of codimension-two embeddings.

%%%%%%%%%%%%%%%%%%%%%%%%%%%%%%%%%%%%%%%%%%%%%%%%%%%%%%%%%%%%
\subsection{Knots and Fundamental Groups}
\label{subsec:knots-fundamental-groups}
%%%%%%%%%%%%%%%%%%%%%%%%%%%%%%%%%%%%%%%%%%%%%%%%%%%%%%%%%%%%

Knot complements provide natural examples of spaces with rich fundamental
groups.

The general pattern is

\[
\boxed{
K
\longrightarrow
S^3\setminus K
\longrightarrow
\pi_1(S^3\setminus K).
}
\]

Thus geometry is converted into topology and then into algebra.

This is one of the clearest examples of the philosophy of algebraic
topology.

%%%%%%%%%%%%%%%%%%%%%%%%%%%%%%%%%%%%%%%%%%%%%%%%%%%%%%%%%%%%
\subsection{Knots and Covering Spaces}
\label{subsec:knots-covering-spaces}
%%%%%%%%%%%%%%%%%%%%%%%%%%%%%%%%%%%%%%%%%%%%%%%%%%%%%%%%%%%%

The abelianization map

\[
\pi_1(X_K)\to\Z
\]

determines an infinite cyclic covering space of the knot exterior.

The homology of this covering space carries an action of the deck
transformation group.

This structure leads naturally to the Alexander module and Alexander
polynomial.

%%%%%%%%%%%%%%%%%%%%%%%%%%%%%%%%%%%%%%%%%%%%%%%%%%%%%%%%%%%%
\subsection{Alexander Module}
\label{subsec:alexander-module}
%%%%%%%%%%%%%%%%%%%%%%%%%%%%%%%%%%%%%%%%%%%%%%%%%%%%%%%%%%%%

Let

\[
\widetilde{X}_K
\]

denote the infinite cyclic cover of the knot exterior.

Its first homology

\[
\boxed{
H_1(\widetilde{X}_K;\Z)
}
\]

becomes a module over the Laurent polynomial ring

\[
\boxed{
\Z[t,t^{-1}].
}
\]

This is called the Alexander module.

The Alexander polynomial can be extracted from a presentation of this
module.

%%%%%%%%%%%%%%%%%%%%%%%%%%%%%%%%%%%%%%%%%%%%%%%%%%%%%%%%%%%%
\subsection{Knots and Homology}
\label{subsec:knots-homology}
%%%%%%%%%%%%%%%%%%%%%%%%%%%%%%%%%%%%%%%%%%%%%%%%%%%%%%%%%%%%

Although knot groups vary greatly,

\[
H_1(X_K;\Z)
\]

is always

\[
\boxed{
\Z.
}
\]

For a knot exterior,

\[
\boxed{
H_0(X_K;\Z)\cong\Z,
}
\]

\[
\boxed{
H_1(X_K;\Z)\cong\Z,
}
\]

and the higher reduced homology is trivial.

Thus ordinary homology alone contains relatively little information about a
single knot.

%%%%%%%%%%%%%%%%%%%%%%%%%%%%%%%%%%%%%%%%%%%%%%%%%%%%%%%%%%%%
\subsection{Alexander Duality}
\label{subsec:knots-alexander-duality}
%%%%%%%%%%%%%%%%%%%%%%%%%%%%%%%%%%%%%%%%%%%%%%%%%%%%%%%%%%%%

Alexander duality relates the topology of a subset of a sphere to the
topology of its complement.

For a sufficiently nice compact subset

\[
A\subset S^n,
\]

one has a relationship of the form

\[
\boxed{
\widetilde{H}_i(S^n\setminus A)
\cong
\widetilde{H}^{\,n-i-1}(A).
}
\]

For

\[
A=S^1\subset S^3,
\]

this explains why knot complements have the same ordinary homology even
though their fundamental groups can be very different.

%%%%%%%%%%%%%%%%%%%%%%%%%%%%%%%%%%%%%%%%%%%%%%%%%%%%%%%%%%%%
\subsection{Knot Theory and Physics}
\label{subsec:knot-theory-physics}
%%%%%%%%%%%%%%%%%%%%%%%%%%%%%%%%%%%%%%%%%%%%%%%%%%%%%%%%%%%%

Knot invariants appear naturally in mathematical physics.

Important connections include

\[
\boxed{
\text{Chern--Simons theory},
}
\]

\[
\boxed{
\text{quantum groups},
}
\]

\[
\boxed{
\text{statistical mechanics},
}
\]

and

\[
\boxed{
\text{topological quantum field theory}.
}
\]

These connections transformed modern knot theory and led to powerful new
invariants.

%%%%%%%%%%%%%%%%%%%%%%%%%%%%%%%%%%%%%%%%%%%%%%%%%%%%%%%%%%%%
\subsection{Knot Theory and Biology}
\label{subsec:knot-theory-biology}
%%%%%%%%%%%%%%%%%%%%%%%%%%%%%%%%%%%%%%%%%%%%%%%%%%%%%%%%%%%%

Long molecular chains can become knotted or linked.

Examples arise in

\[
\boxed{
\text{DNA},
\quad
\text{proteins},
\quad
\text{polymers}.
}
\]

Topological methods can help describe changes caused by cutting,
reconnecting, twisting, or passing molecular strands through one another.

%%%%%%%%%%%%%%%%%%%%%%%%%%%%%%%%%%%%%%%%%%%%%%%%%%%%%%%%%%%%
\subsection{Knot Theory and Chemistry}
\label{subsec:knot-theory-chemistry}
%%%%%%%%%%%%%%%%%%%%%%%%%%%%%%%%%%%%%%%%%%%%%%%%%%%%%%%%%%%%

Chemists have constructed molecular knots and links.

Their topology can influence

\[
\boxed{
\text{molecular stability},
\quad
\text{self-assembly},
\quad
\text{mechanical properties}.
}
\]

Thus knot theory provides a language for molecular architecture.

%%%%%%%%%%%%%%%%%%%%%%%%%%%%%%%%%%%%%%%%%%%%%%%%%%%%%%%%%%%%
\subsection{Knot Theory and Fluid Dynamics}
\label{subsec:knot-theory-fluid-dynamics}
%%%%%%%%%%%%%%%%%%%%%%%%%%%%%%%%%%%%%%%%%%%%%%%%%%%%%%%%%%%%

Vortex lines in fluids may form knots and links.

Their topology can influence the evolution and interaction of vortex
structures.

Related ideas appear in

\[
\boxed{
\text{fluid mechanics},
\quad
\text{plasma physics},
\quad
\text{magnetic fields}.
}
\]

%%%%%%%%%%%%%%%%%%%%%%%%%%%%%%%%%%%%%%%%%%%%%%%%%%%%%%%%%%%%
\subsection{Knot Theory and Computation}
\label{subsec:knot-theory-computation}
%%%%%%%%%%%%%%%%%%%%%%%%%%%%%%%%%%%%%%%%%%%%%%%%%%%%%%%%%%%%

Computational knot theory studies algorithms for tasks such as

\[
\boxed{
\text{unknot recognition},
}
\]

\[
\boxed{
\text{invariant computation},
}
\]

\[
\boxed{
\text{diagram simplification},
}
\]

and

\[
\boxed{
\text{three-manifold analysis}.
}
\]

A diagram may be easy to draw but difficult to simplify.

Thus knot theory contains significant algorithmic complexity.

%%%%%%%%%%%%%%%%%%%%%%%%%%%%%%%%%%%%%%%%%%%%%%%%%%%%%%%%%%%%
\subsection{The Unknot Recognition Problem}
\label{subsec:unknot-recognition}
%%%%%%%%%%%%%%%%%%%%%%%%%%%%%%%%%%%%%%%%%%%%%%%%%%%%%%%%%%%%

Given a knot diagram, determine whether it represents the unknot.

This is called the \emph{unknot recognition problem}.

A diagram may contain many crossings even when the underlying knot is
trivial.

Therefore visual inspection alone is unreliable.

Mathematical algorithms use structures such as normal surfaces and
\(3\)-manifold topology.

%%%%%%%%%%%%%%%%%%%%%%%%%%%%%%%%%%%%%%%%%%%%%%%%%%%%%%%%%%%%
\subsection{A Hierarchy of Knot Information}
\label{subsec:hierarchy-knot-information}
%%%%%%%%%%%%%%%%%%%%%%%%%%%%%%%%%%%%%%%%%%%%%%%%%%%%%%%%%%%%

Knot theory provides many levels of information:

\[
\boxed{
\text{diagram}
}
\]

\[
\Downarrow
\]

\[
\boxed{
\text{elementary invariants}
}
\]

\[
\Downarrow
\]

\[
\boxed{
\text{polynomial invariants}
}
\]

\[
\Downarrow
\]

\[
\boxed{
\text{knot group and covering spaces}
}
\]

\[
\Downarrow
\]

\[
\boxed{
\text{knot complement geometry}
}
\]

\[
\Downarrow
\]

\[
\boxed{
\text{homological and four-dimensional invariants}.
}
\]

No single invariant introduced here completely replaces the others.

%%%%%%%%%%%%%%%%%%%%%%%%%%%%%%%%%%%%%%%%%%%%%%%%%%%%%%%%%%%%
\subsection{Worked Example: Alexander Polynomial}
\label{subsec:worked-alexander-polynomial}
%%%%%%%%%%%%%%%%%%%%%%%%%%%%%%%%%%%%%%%%%%%%%%%%%%%%%%%%%%%%

\begin{workedexamplebox}{Trefoil Determinant}

Suppose

\[
\Delta_K(t)=t^2-t+1.
\]

The knot determinant is

\[
\det(K)
=
|\Delta_K(-1)|.
\]

Substitute

\[
t=-1:
\]

\[
\Delta_K(-1)
=
(-1)^2-(-1)+1.
\]

Thus

\[
\Delta_K(-1)
=
1+1+1
=
3.
\]

\begin{answerbox}
\[
\boxed{
\det(K)=3.
}
\]
\end{answerbox}
\end{workedexamplebox}

%%%%%%%%%%%%%%%%%%%%%%%%%%%%%%%%%%%%%%%%%%%%%%%%%%%%%%%%%%%%
\subsection{Worked Example: Torus Knot Genus}
\label{subsec:worked-torus-knot-genus}
%%%%%%%%%%%%%%%%%%%%%%%%%%%%%%%%%%%%%%%%%%%%%%%%%%%%%%%%%%%%

\begin{workedexamplebox}{Genus of \(T(3,4)\)}

Since

\[
\gcd(3,4)=1,
\]

the curve

\[
T(3,4)
\]

is a knot.

The genus formula is

\[
g(T(p,q))
=
\frac{(p-1)(q-1)}{2}.
\]

Therefore

\[
g(T(3,4))
=
\frac{(3-1)(4-1)}{2}.
\]

Hence

\[
g(T(3,4))
=
\frac{2\cdot3}{2}
=
3.
\]

\begin{answerbox}
\[
\boxed{
g(T(3,4))=3.
}
\]
\end{answerbox}
\end{workedexamplebox}

%%%%%%%%%%%%%%%%%%%%%%%%%%%%%%%%%%%%%%%%%%%%%%%%%%%%%%%%%%%%
\subsection{Worked Example: Torus Knot Crossing Number}
\label{subsec:worked-torus-crossing-number}
%%%%%%%%%%%%%%%%%%%%%%%%%%%%%%%%%%%%%%%%%%%%%%%%%%%%%%%%%%%%

\begin{workedexamplebox}{Crossing Number of \(T(3,5)\)}

For positive coprime \(p,q\),

\[
c(T(p,q))
=
\min\{(p-1)q,(q-1)p\}.
\]

For

\[
T(3,5),
\]

we obtain

\[
(p-1)q
=
2(5)
=
10,
\]

while

\[
(q-1)p
=
4(3)
=
12.
\]

Therefore

\[
\begin{answerbox}
\[
\boxed{
c(T(3,5))=10.
}
\]
\end{answerbox}
\]
\end{workedexamplebox}

%%%%%%%%%%%%%%%%%%%%%%%%%%%%%%%%%%%%%%%%%%%%%%%%%%%%%%%%%%%%
\subsection{Common Errors}
\label{subsec:knot-theory-common-errors}
%%%%%%%%%%%%%%%%%%%%%%%%%%%%%%%%%%%%%%%%%%%%%%%%%%%%%%%%%%%%

\subsubsection{Confusing a Knot with a Knot Diagram}

A knot lives in three-dimensional space.

A diagram is a two-dimensional representation of the knot.

%%%%%%%%%%%%%%%%%%%%%%%%%%%%%%%%%%%%%%%%%%%%%%%%%%%%%%%%%%%%
\subsubsection{Assuming Different Diagrams Represent Different Knots}

Reidemeister moves can dramatically change a diagram without changing the
underlying knot.

%%%%%%%%%%%%%%%%%%%%%%%%%%%%%%%%%%%%%%%%%%%%%%%%%%%%%%%%%%%%
\subsubsection{Treating Crossings as Physical Intersections}

The strands do not actually intersect in three-dimensional space.

One strand passes over the other.

%%%%%%%%%%%%%%%%%%%%%%%%%%%%%%%%%%%%%%%%%%%%%%%%%%%%%%%%%%%%
\subsubsection{Assuming Crossing Number Is the Number of Crossings in Any Diagram}

The crossing number is the \emph{minimum} over all diagrams:

\[
\boxed{
c(K)
=
\min_D c(D).
}
\]

%%%%%%%%%%%%%%%%%%%%%%%%%%%%%%%%%%%%%%%%%%%%%%%%%%%%%%%%%%%%
\subsubsection{Assuming Zero Linking Number Means Unlinked}

A link can have

\[
\operatorname{lk}=0
\]

and still be nontrivial.

%%%%%%%%%%%%%%%%%%%%%%%%%%%%%%%%%%%%%%%%%%%%%%%%%%%%%%%%%%%%
\subsubsection{Forgetting Orientation in Linking Number}

Changing the orientation of one component changes the sign of the linking
number.

%%%%%%%%%%%%%%%%%%%%%%%%%%%%%%%%%%%%%%%%%%%%%%%%%%%%%%%%%%%%
\subsubsection{Assuming the Knot Group Completely Classifies Knots by Abstract Group Alone}

The knot group is powerful, but peripheral information is important when
recovering the knot from its complement.

One should distinguish the abstract group from the full geometric structure
of the knot exterior.

%%%%%%%%%%%%%%%%%%%%%%%%%%%%%%%%%%%%%%%%%%%%%%%%%%%%%%%%%%%%
\subsubsection{Assuming an Alexander Polynomial of One Means Unknot}

Nontrivial knots can satisfy

\[
\boxed{
\Delta_K(t)=1.
}
\]

%%%%%%%%%%%%%%%%%%%%%%%%%%%%%%%%%%%%%%%%%%%%%%%%%%%%%%%%%%%%
\subsubsection{Ignoring Polynomial Normalization}

The Alexander polynomial is naturally defined only up to

\[
\boxed{
\pm t^m.
}
\]

Different books may therefore display equivalent forms.

%%%%%%%%%%%%%%%%%%%%%%%%%%%%%%%%%%%%%%%%%%%%%%%%%%%%%%%%%%%%
\subsubsection{Confusing Seifert Surface Genus with Knot Genus}

A particular Seifert surface may not have minimum genus.

The knot genus is the minimum over all Seifert surfaces.

%%%%%%%%%%%%%%%%%%%%%%%%%%%%%%%%%%%%%%%%%%%%%%%%%%%%%%%%%%%%
\subsubsection{Assuming the Seifert Algorithm Always Produces a Minimum-Genus Surface}

The algorithm constructs a Seifert surface, but it need not minimize genus.

Thus it generally gives an upper bound.

%%%%%%%%%%%%%%%%%%%%%%%%%%%%%%%%%%%%%%%%%%%%%%%%%%%%%%%%%%%%
\subsubsection{Assuming Every \texorpdfstring{\(T(p,q)\)}{T(p,q)} Is a Knot}

If

\[
\gcd(p,q)>1,
\]

then \(T(p,q)\) has multiple components and is a torus link.

%%%%%%%%%%%%%%%%%%%%%%%%%%%%%%%%%%%%%%%%%%%%%%%%%%%%%%%%%%%%
\subsubsection{Confusing Writhe with a Knot Invariant}

Writhe changes under Reidemeister move I.

Therefore it depends on the diagram.

%%%%%%%%%%%%%%%%%%%%%%%%%%%%%%%%%%%%%%%%%%%%%%%%%%%%%%%%%%%%
\subsubsection{Assuming Every Knot Is Hyperbolic}

Torus knots and satellite knots are not hyperbolic knots.

%%%%%%%%%%%%%%%%%%%%%%%%%%%%%%%%%%%%%%%%%%%%%%%%%%%%%%%%%%%%
\subsubsection{Confusing Slice with Unknotted}

A slice knot may be nontrivial in \(S^3\).

Slice means it bounds an embedded disk in \(D^4\).

%%%%%%%%%%%%%%%%%%%%%%%%%%%%%%%%%%%%%%%%%%%%%%%%%%%%%%%%%%%%
\subsection{Applications and Connections}
\label{subsec:knot-theory-applications}
%%%%%%%%%%%%%%%%%%%%%%%%%%%%%%%%%%%%%%%%%%%%%%%%%%%%%%%%%%%%

\begin{applicationbox}
\textbf{Geometric Topology.}
Knots are embeddings of circles in \(3\)-manifolds and provide fundamental
examples of embedding problems.

\medskip

\textbf{Algebraic Topology.}
Knot groups, covering spaces, homology, and Alexander modules translate
geometric knotting into algebraic information.

\medskip

\textbf{Three-Manifold Topology.}
Knot exteriors are \(3\)-manifolds with torus boundary. Dehn filling and
geometrization connect knot theory directly with manifold classification.

\medskip

\textbf{Hyperbolic Geometry.}
Many knot complements admit complete finite-volume hyperbolic structures.

\medskip

\textbf{Group Theory.}
Wirtinger presentations and braid groups provide natural nonabelian groups
arising from knots.

\medskip

\textbf{Linear Algebra.}
Seifert matrices encode linking information and produce polynomial
invariants through determinants.

\medskip

\textbf{Four-Dimensional Topology.}
Concordance and slice knots study how knots bound surfaces in the \(4\)-ball.

\medskip

\textbf{Mathematical Physics.}
Jones-type invariants connect knots with quantum groups, Chern--Simons
theory, statistical mechanics, and quantum topology.

\medskip

\textbf{Biology and Chemistry.}
DNA, proteins, polymers, and molecular structures can exhibit knotting and
linking.

\medskip

\textbf{Computation.}
Knot recognition and invariant computation provide important algorithmic
problems in topology.
\end{applicationbox}

%%%%%%%%%%%%%%%%%%%%%%%%%%%%%%%%%%%%%%%%%%%%%%%%%%%%%%%%%%%%
\subsection{Exercises}
\label{subsec:knot-theory-exercises}
%%%%%%%%%%%%%%%%%%%%%%%%%%%%%%%%%%%%%%%%%%%%%%%%%%%%%%%%%%%%

\begin{exercise}[1]
Define a knot mathematically.
\end{exercise}

\begin{exercise}[1]
What is the unknot?
\end{exercise}

\begin{exercise}[1]
What is a knot diagram?
\end{exercise}

\begin{exercise}[1]
What is ambient isotopy?
\end{exercise}

\begin{exercise}[1]
State the three Reidemeister moves.
\end{exercise}

\begin{exercise}[1]
Define crossing number.
\end{exercise}

\begin{exercise}[1]
Define a prime knot.
\end{exercise}

\begin{exercise}[1]
Define a link.
\end{exercise}

\begin{exercise}[1]
Define linking number.
\end{exercise}

\begin{exercise}[1]
Define the knot group.
\end{exercise}

\begin{exercise}[1]
What is a Seifert surface?
\end{exercise}

\begin{exercise}[1]
Define knot genus.
\end{exercise}

\begin{exercise}[1]
What is a braid?
\end{exercise}

\begin{exercise}[1]
What is a torus knot?
\end{exercise}

\begin{exercise}[1]
Define a slice knot.
\end{exercise}

\begin{exercise}[2]
Explain why the number of crossings in a particular diagram need not equal
the crossing number of the knot.
\end{exercise}

\begin{exercise}[2]
Which Reidemeister move changes the crossing number of a diagram by one?
\end{exercise}

\begin{exercise}[2]
Which Reidemeister move introduces or removes two crossings?
\end{exercise}

\begin{exercise}[2]
Which Reidemeister move preserves the number of crossings?
\end{exercise}

\begin{exercise}[2]
Suppose mutual crossings of two oriented link components have signs

\[
+1,+1,-1,+1,-1,+1.
\]

Compute their linking number.
\end{exercise}

\begin{exercise}[2]
Explain why nonzero linking number proves that a two-component link is
nontrivial.
\end{exercise}

\begin{exercise}[2]
Explain why zero linking number does not prove that a link is trivial.
\end{exercise}

\begin{exercise}[2]
What is the knot group of the unknot?
\end{exercise}

\begin{exercise}[2]
Explain why a knot with nonabelian knot group cannot be the unknot.
\end{exercise}

\begin{exercise}[2]
State the Fox \(n\)-coloring relation.
\end{exercise}

\begin{exercise}[2]
Explain why three-colorability distinguishes the trefoil from the unknot.
\end{exercise}

\begin{exercise}[2]
If

\[
\Delta_K(t)=t^2-3t+1,
\]

calculate \(\det(K)\).
\end{exercise}

\begin{exercise}[2]
If

\[
\Delta_K(t)=1,
\]

can you conclude that \(K\) is the unknot? Explain.
\end{exercise}

\begin{exercise}[2]
Suppose a Seifert algorithm produces

\[
c=9
\]

crossings and

\[
s=6
\]

Seifert circles. Determine the genus of the resulting surface.
\end{exercise}

\begin{exercise}[2]
Compute

\[
g(T(2,7)).
\]
\end{exercise}

\begin{exercise}[2]
Compute

\[
g(T(4,5)).
\]
\end{exercise}

\begin{exercise}[2]
Compute

\[
c(T(2,9)).
\]
\end{exercise}

\begin{exercise}[2]
Determine whether \(T(4,6)\) is a knot or a link.
\end{exercise}

\begin{exercise}[2]
Determine the number of components of \(T(4,6)\).
\end{exercise}

\begin{exercise}[3]
Use Reidemeister moves to simplify several diagrams of the unknot.
\end{exercise}

\begin{exercise}[3]
Construct a nontrivial Fox \(3\)-coloring of a trefoil diagram.
\end{exercise}

\begin{exercise}[3]
Derive a Wirtinger presentation for the trefoil knot group.
\end{exercise}

\begin{exercise}[3]
Show that the trefoil group

\[
\langle a,b\mid aba=bab\rangle
\]

is nonabelian.
\end{exercise}

\begin{exercise}[3]
Show that the abelianization of the trefoil knot group is \(\Z\).
\end{exercise}

\begin{exercise}[3]
Explain why the abelianization of every knot group is \(\Z\).
\end{exercise}

\begin{exercise}[3]
Use the Seifert algorithm on an oriented trefoil diagram and determine the
genus of the resulting surface.
\end{exercise}

\begin{exercise}[3]
Show that

\[
g(K_1\#K_2)
=
g(K_1)+g(K_2).
\]
\end{exercise}

\begin{exercise}[3]
Show that the Alexander polynomial is multiplicative under connected sum,
up to units.
\end{exercise}

\begin{exercise}[3]
Verify the braid relation

\[
\sigma_1\sigma_2\sigma_1
=
\sigma_2\sigma_1\sigma_2
\]

geometrically using braid diagrams.
\end{exercise}

\begin{exercise}[3]
Write the trefoil as the closure of a braid.
\end{exercise}

\begin{exercise}[3]
Show that the closure of

\[
\sigma_1^3\in B_2
\]

is a trefoil knot.
\end{exercise}

\begin{exercise}[3]
Explain why \(T(p,q)\) has \(\gcd(p,q)\) components.
\end{exercise}

\begin{exercise}[3]
Compute the crossing number and genus of

\[
T(5,7).
\]
\end{exercise}

\begin{exercise}[3]
Explain why writhe fails to be invariant under all Reidemeister moves.
\end{exercise}

\begin{exercise}[3]
Describe how the Kauffman bracket can be normalized using writhe.
\end{exercise}

\begin{exercise}[3]
Explain why the knot exterior has torus boundary.
\end{exercise}

\begin{exercise}[3]
Describe the meridian and preferred longitude of a knot.
\end{exercise}

\begin{exercise}[3]
Explain geometrically how Dehn surgery modifies a knot exterior.
\end{exercise}

\begin{exercise}[3]
Explain why a hyperbolic knot has a well-defined hyperbolic volume.
\end{exercise}

\begin{exercise}[3]
Compare torus, satellite, and hyperbolic knots.
\end{exercise}

\begin{exercise}[4]
Study a proof of the Reidemeister Theorem.
\end{exercise}

\begin{exercise}[4]
Research the prime decomposition theorem for knots.
\end{exercise}

\begin{exercise}[4]
Study the Wirtinger presentation using the Seifert--van Kampen Theorem.
\end{exercise}

\begin{exercise}[4]
Construct the Alexander polynomial from a presentation of the Alexander
module.
\end{exercise}

\begin{exercise}[4]
Compute an Alexander polynomial using a Seifert matrix.
\end{exercise}

\begin{exercise}[4]
Study the proof of Alexander's Theorem that every oriented link is a closed
braid.
\end{exercise}

\begin{exercise}[4]
Research Markov's Theorem and explain why it is the braid analogue of the
Reidemeister Theorem.
\end{exercise}

\begin{exercise}[4]
Study the Kauffman bracket construction of the Jones polynomial.
\end{exercise}

\begin{exercise}[4]
Research the HOMFLY-PT polynomial and its specializations.
\end{exercise}

\begin{exercise}[4]
Study the figure-eight knot complement as a hyperbolic \(3\)-manifold.
\end{exercise}

\begin{exercise}[4]
Research the decomposition of prime knots into torus, satellite, and
hyperbolic types.
\end{exercise}

\begin{exercise}[4]
Study the theorem that knots are determined by their complements.
\end{exercise}

\begin{exercise}[4]
Research Dehn surgery and the Lickorish--Wallace theorem.
\end{exercise}

\begin{exercise}[4]
Study the smooth knot concordance group.
\end{exercise}

\begin{exercise}[4]
Find an example of a nontrivial slice knot and explain why being slice does
not imply being the unknot.
\end{exercise}

\begin{exercise}[4]
Research Khovanov homology as a categorification of the Jones polynomial.
\end{exercise}

\begin{exercise}[4]
Research knot Floer homology and its relationship with knot genus.
\end{exercise}

\begin{exercise}[4]
Study higher-dimensional knots

\[
S^n\hookrightarrow S^{n+2}.
\]
\end{exercise}

%%%%%%%%%%%%%%%%%%%%%%%%%%%%%%%%%%%%%%%%%%%%%%%%%%%%%%%%%%%%
\subsection{Conceptual Summary}
\label{subsec:knot-theory-summary}
%%%%%%%%%%%%%%%%%%%%%%%%%%%%%%%%%%%%%%%%%%%%%%%%%%%%%%%%%%%%

A classical knot is an embedding

\[
\boxed{
S^1\hookrightarrow S^3.
}
\]

Two knots are equivalent when they are related by ambient isotopy.

Knot diagrams represent three-dimensional knots in two dimensions.

The Reidemeister moves

\[
\boxed{
R_1,\quad R_2,\quad R_3
}
\]

completely characterize diagrammatic equivalence for tame knots.

Knot invariants provide tools for distinguishing knot types.

Elementary invariants include

\[
\boxed{
\text{crossing number},
\quad
\text{colorability},
\quad
\text{genus}.
}
\]

Algebraic invariants include

\[
\boxed{
\pi_1(S^3\setminus K)
}
\]

and polynomial invariants such as

\[
\boxed{
\Delta_K(t),
\quad
V_K(t).
}
\]

Links generalize knots to multiple components.

For oriented links, linking number measures a basic form of pairwise
linking:

\[
\boxed{
\operatorname{lk}(L_1,L_2)
=
\frac12
\sum_c\operatorname{sgn}(c).
}
\]

Seifert surfaces connect knots with surface topology.

The knot genus is

\[
\boxed{
g(K)
=
\min\{
g(F):
\partial F=K
\}.
}
\]

Braids convert knot theory into group theory.

Every oriented link can be represented as the closure of a braid.

Torus knots provide an important explicit family:

\[
\boxed{
T(p,q).
}
\]

For coprime positive \(p,q\),

\[
\boxed{
g(T(p,q))
=
\frac{(p-1)(q-1)}{2}.
}
\]

Knot complements connect knot theory with \(3\)-manifold topology.

Many prime knots are hyperbolic, making hyperbolic volume a knot invariant.

Dehn surgery transforms knot exteriors into new \(3\)-manifolds.

Concordance and slice knots extend the theory into four dimensions.

Thus knot theory forms a major intersection of

\[
\boxed{
\text{topology},
\quad
\text{geometry},
\quad
\text{algebra},
\quad
\text{combinatorics},
\quad
\text{physics}.
}
\]

%%%%%%%%%%%%%%%%%%%%%%%%%%%%%%%%%%%%%%%%%%%%%%%%%%%%%%%%%%%%
\subsection{Section Connections}
\label{subsec:knot-theory-connections}
%%%%%%%%%%%%%%%%%%%%%%%%%%%%%%%%%%%%%%%%%%%%%%%%%%%%%%%%%%%%

\chapterconnection{
Knot Theory develops one of the central embedding problems of topology:
understanding copies of \(S^1\) inside \(S^3\). Point-Set Topology provides
the language of embeddings and homeomorphisms. Algebraic Topology supplies
fundamental groups, covering spaces, homology, Alexander duality, and the
Seifert--van Kampen Theorem. Geometric Topology provides ambient isotopy,
three-manifolds, connected sums, Dehn surgery, and hyperbolic structures.
Manifold Theory explains knot exteriors as \(3\)-manifolds with torus
boundary. Group Theory appears through knot groups and braid groups. Linear
Algebra appears through Seifert matrices. Hyperbolic Geometry describes many
knot complements. Four-dimensional topology enters through concordance and
slice knots. The next section develops Low-Dimensional Topology, where knots,
surfaces, \(3\)-manifolds, and \(4\)-manifolds become part of a broader
dimension-specific theory.
}

%%%%%%%%%%%%%%%%%%%%%%%%%%%%%%%%%%%%%%%%%%%%%%%%%%%%%%%%%%%%
% SECTION SOURCE TRACKING
%%%%%%%%%%%%%%%%%%%%%%%%%%%%%%%%%%%%%%%%%%%%%%%%%%%%%%%%%%%%

%------------------------------------------------------------
% 4.6 Knot Theory
%------------------------------------------------------------
%
% Source basis:
%   - Standard introductory knot theory, algebraic topology,
%     braid theory, and 3-manifold topology.
%
% Used for:
%   - knots as embeddings S^1 -> S^3;
%   - ambient isotopy;
%   - unknot and nontrivial knots;
%   - knot diagrams;
%   - regular projections;
%   - Reidemeister moves and theorem;
%   - knot invariants;
%   - crossing number;
%   - knot tables;
%   - orientations;
%   - mirror images and chirality;
%   - connected sums;
%   - prime decomposition;
%   - links and unlink;
%   - Hopf link;
%   - crossing signs;
%   - linking number;
%   - knot complements and exteriors;
%   - knot groups;
%   - Wirtinger presentations;
%   - trefoil group;
%   - abelianization;
%   - Fox colorings;
%   - Alexander polynomial;
%   - Jones polynomial;
%   - skein relations;
%   - HOMFLY-PT polynomial;
%   - Seifert surfaces;
%   - knot genus;
%   - Seifert algorithm;
%   - Seifert matrices;
%   - braid groups;
%   - Alexander's braid theorem;
%   - Markov moves;
%   - torus knots;
%   - alternating knots;
%   - determinant;
%   - writhe;
%   - Kauffman bracket;
%   - hyperbolic knots;
%   - hyperbolic volume;
%   - torus/satellite/hyperbolic decomposition;
%   - satellite knots;
%   - Dehn surgery;
%   - meridian and longitude;
%   - concordance;
%   - slice knots;
%   - concordance group;
%   - higher-dimensional knots;
%   - Alexander module;
%   - Alexander duality;
%   - applications.
%
% Notes:
%   - Fundamental groups, covering spaces, homology, and
%     Alexander duality are developed canonically in
%     Algebraic Topology.
%   - Ambient isotopy, prime decomposition, hyperbolic
%     structures, and Dehn surgery connect to Geometric
%     Topology.
%   - Manifold foundations and torus boundaries are
%     developed in Manifold Theory.
%   - Hyperbolic geometry is developed canonically in
%     Non-Euclidean Geometry and Differential Geometry.
%   - Concordance and slice knots provide a bridge to
%     4-manifold topology.
%   - Khovanov homology, Floer homology, quantum topology,
%     and advanced concordance invariants are introduced
%     only structurally here.
%
%%%%%%%%%%%%%%%%%%%%%%%%%%%%%%%%%%%%%%%%%%%%%%%%%%%%%%%%%%%%